% Solutions to Homework 8 — Math Physics Methods (数学物理方法)

\PassOptionsToPackage{quiet}{xeCJK}
\documentclass[11pt]{article}

\usepackage[margin=1in]{geometry}
\usepackage{amsmath,amsthm,amssymb,mathtools}
\usepackage{ctex}
\usepackage{fancyhdr}
\usepackage{fontspec}
\usepackage{fourier-orns}

\theoremstyle{remark}
\newtheorem{problem}{}
\theoremstyle{definition}
\newtheorem{solution}{解}

\renewcommand{\headrule}{%
\vspace{-8pt}\hrulefill
\raisebox{-2.0pt}{\quad\decofourleft\decotwo\decofourright\quad}\hrulefill}

\setlength{\headheight}{15.11996pt}
\setlength{\topmargin}{-3.11996pt}

\newcommand{\duedate}{\zhdate{2026/12/01}}
\pagestyle{fancy}

\fancyhead{}
\fancyhead[L]{数学物理方法}
\zhnumsetup{style={Traditional,Financial}}
\fancyhead[R]{课后习题捌 \quad 参考解答}

\fancyfoot{}
\fancyfoot[C]{\thepage}
\fancyfoot[L]{截止日期\duedate}
\fancyfoot[R]{任课教师:罗凯}

\begin{document}
\renewcommand{\labelenumi}{(\arabic{enumi})}
\renewcommand{\labelenumii}{(\arabic{enumi}.\arabic{enumii})}

%==========================================================================
\begin{problem}
  用爱因斯坦求和约定改写下列方程. 指标取 $1,2,3$，重复指标表示求和，式中不再出现 $\sum$. 叉积和旋度用 $\varepsilon_{ijk}$.
  \begin{enumerate}
    \item 连续性方程 $\partial_t\rho+\nabla\cdot\mathbf{J}=0$;
    \item 泊松方程 $\nabla^2\varphi=-\rho/\varepsilon_0$;
    \item 真空中的麦克斯韦方程组
      \[
        \begin{aligned}
          \nabla\cdot\mathbf{E}&=\rho/\varepsilon_0, &
          \nabla\cdot\mathbf{B}&=0,\\
          \nabla\times\mathbf{E}&=-\partial_t\mathbf{B}, &
          \nabla\times\mathbf{B}
          &=\mu_0\mathbf{J}+\mu_0\varepsilon_0\partial_t\mathbf{E};
        \end{aligned}
      \]
    \item 无源波动方程
      $\nabla^2\mathbf{E}-c^{-2}\partial_{tt}\mathbf{E}=\mathbf{0}$;
    \item 洛伦兹力
      $\mathbf{F}=q\bigl(\mathbf{E}+\mathbf{v}\times\mathbf{B}\bigr)$.
  \end{enumerate}
\end{problem}

\begin{solution}
直角坐标里 $\nabla\cdot\mathbf{J}=\partial_i J_i$，$\nabla^2=\partial_i\partial_i$，
$(\nabla\times\mathbf{A})_i=\varepsilon_{ijk}\partial_j A_k$，
$(\mathbf{v}\times\mathbf{B})_i=\varepsilon_{ijk}v_j B_k$。矢量方程写成第 $i$ 个分量。

\begin{enumerate}
  \item $\partial_t\rho+\partial_i J_i=0$.
  \item $\partial_i\partial_i\varphi=-\rho/\varepsilon_0$.
  \item
    \[
      \begin{aligned}
        \partial_i E_i&=\rho/\varepsilon_0, &
        \partial_i B_i&=0,\\
        \varepsilon_{ijk}\partial_j E_k&=-\partial_t B_i, &
        \varepsilon_{ijk}\partial_j B_k
        &=\mu_0 J_i+\mu_0\varepsilon_0\partial_t E_i.
      \end{aligned}
    \]
  \item 直角坐标中矢量拉普拉斯按分量作用，
    $(\nabla^2\mathbf{E})_i=\partial_j\partial_j E_i$，故
    \[
      \partial_j\partial_j E_i-c^{-2}\partial_{tt}E_i=0.
    \]
  \item $F_i=q\bigl(E_i+\varepsilon_{ijk}v_j B_k\bigr)$.
\end{enumerate}
\end{solution}

%==========================================================================
\begin{problem}
  按爱因斯坦约定去掉求和号，再缩并成更短的形式.
  \begin{enumerate}
    \item $\displaystyle\sum_{j=1}^{3}\delta_{ij}\,\partial_j\varphi$;
    \item $\displaystyle\sum_{i=1}^{3}\sum_{j=1}^{3}\delta_{ij}\,\partial_i A_j$;
    \item $\displaystyle\sum_{j=1}^{3}\sum_{k=1}^{3}\varepsilon_{ijk}\,\delta_{jk}$;
    \item 不可压缩流动
      \[
        \partial_t u_i
        +\sum_{j=1}^{3}u_j\partial_j u_i
        =-\frac{1}{\rho}\sum_{j=1}^{3}\delta_{ij}\,\partial_j p
         +\nu\sum_{j=1}^{3}\partial_j\partial_j u_i,
        \qquad
        \sum_{i=1}^{3}\partial_i u_i=0.
      \]
  \end{enumerate}
\end{problem}

\begin{solution}
\begin{enumerate}
  \item $\delta_{ij}$ 只在 $j=i$ 时为 $1$，故
    \[
      \delta_{ij}\,\partial_j\varphi=\partial_i\varphi.
    \]
    这是 $(\nabla\varphi)_i$。自由指标是 $i$。
  \item $i,j$ 都是哑指标。$\delta_{ij}$ 令 $j=i$，
    \[
      \delta_{ij}\,\partial_i A_j=\partial_i A_i=\nabla\cdot\mathbf{A}.
    \]
    结果是标量，没有自由指标。
  \item $\delta_{jk}$ 令 $k=j$，剩下 $\varepsilon_{ijj}$。两个指标相同，$\varepsilon_{ijj}=0$，所以
    \[
      \varepsilon_{ijk}\,\delta_{jk}=0.
    \]
    自由指标仍是 $i$，三个分量都是零。
  \item $\delta_{ij}\partial_j p=\partial_i p$，其余只是去掉求和号：
    \[
      \partial_t u_i+u_j\partial_j u_i
      =-\frac{1}{\rho}\partial_i p+\nu\,\partial_j\partial_j u_i,
      \qquad
      \partial_i u_i=0.
    \]
\end{enumerate}
\end{solution}

%==========================================================================
\begin{problem}
  用指标记号证明下列乘积法则. $f$ 是标量场, $\mathbf{A}$、$\mathbf{B}$ 是矢量场, 分量在直角坐标中足够光滑.
  \begin{enumerate}
    \item $\displaystyle
      \nabla\cdot(f\mathbf{A})
      =f\,\nabla\cdot\mathbf{A}+\mathbf{A}\cdot\nabla f$;
    \item $\displaystyle
      \nabla\times(f\mathbf{A})
      =f\,\nabla\times\mathbf{A}+\nabla f\times\mathbf{A}$;
    \item $\displaystyle
      \nabla\cdot(\mathbf{A}\times\mathbf{B})
      =\mathbf{B}\cdot(\nabla\times\mathbf{A})
       -\mathbf{A}\cdot(\nabla\times\mathbf{B})$.
  \end{enumerate}
\end{problem}

\begin{solution}
\begin{enumerate}
  \item
    \[
      \nabla\cdot(f\mathbf{A})
      =\partial_i(f A_i)
      =(\partial_i f)A_i+f\,\partial_i A_i
      =\mathbf{A}\cdot\nabla f+f\,\nabla\cdot\mathbf{A}.
    \]
  \item 第 $i$ 个分量
    \begin{align*}
      \bigl(\nabla\times(f\mathbf{A})\bigr)_i
      &=\varepsilon_{ijk}\partial_j(f A_k)
       =\varepsilon_{ijk}(\partial_j f)A_k
        +f\,\varepsilon_{ijk}\partial_j A_k \\
      &=(\nabla f\times\mathbf{A})_i
        +f\,(\nabla\times\mathbf{A})_i.
    \end{align*}
  \item $(\mathbf{A}\times\mathbf{B})_i=\varepsilon_{ijk}A_j B_k$，故
    \[
      \nabla\cdot(\mathbf{A}\times\mathbf{B})
      =\varepsilon_{ijk}\partial_i(A_j B_k)
      =\varepsilon_{ijk}(\partial_i A_j)B_k
       +\varepsilon_{ijk}A_j(\partial_i B_k).
    \]
    第一项：$(i,j,k)\mapsto(k,i,j)$ 是偶置换，$\varepsilon_{kij}=\varepsilon_{ijk}$，而
    $(\nabla\times\mathbf{A})_k=\varepsilon_{kij}\partial_i A_j$，所以
    \[
      \varepsilon_{ijk}(\partial_i A_j)B_k
      =B_k(\nabla\times\mathbf{A})_k
      =\mathbf{B}\cdot(\nabla\times\mathbf{A}).
    \]
    第二项：对换前两个指标，$\varepsilon_{jik}=-\varepsilon_{ijk}$，而
    $(\nabla\times\mathbf{B})_j=\varepsilon_{jik}\partial_i B_k$，所以
    \[
      \varepsilon_{ijk}A_j(\partial_i B_k)
      =-A_j(\nabla\times\mathbf{B})_j
      =-\mathbf{A}\cdot(\nabla\times\mathbf{B}).
    \]
    两项相加即得。
\end{enumerate}
\end{solution}

%==========================================================================
\begin{problem}
  \begin{enumerate}
    \item 设 $\mathbf{A}=x\,\mathbf{e}_x+y\,\mathbf{e}_y+z\,\mathbf{e}_z$，$V$ 是球 $x^2+y^2+z^2\le R^2$，$S=\partial V$ 取外法向. 分别计算散度定理的两边，验证
      \[
        \oint_S\mathbf{A}\cdot\mathrm{d}\mathbf{S}
        =\int_V(\nabla\cdot\mathbf{A})\,\mathrm{d}V.
      \]
    \item 设 $\mathbf{A}=y\,\mathbf{e}_x+z\,\mathbf{e}_y+x\,\mathbf{e}_z$. $S$ 是正方形 $0\le x\le a$，$0\le y\le a$，$z=0$，法向取 $\mathbf{e}_z$，边界 $L$ 按右手定则定向. 分别计算斯托克斯定理的两边，验证
      \[
        \oint_L\mathbf{A}\cdot\mathrm{d}\boldsymbol{\ell}
        =\int_S(\nabla\times\mathbf{A})\cdot\mathrm{d}\mathbf{S}.
      \]
  \end{enumerate}
\end{problem}

\begin{solution}
\begin{enumerate}
  \item $\nabla\cdot\mathbf{A}=\partial_x x+\partial_y y+\partial_z z=3$。球的体积是 $\tfrac43\pi R^3$，
    \[
      \int_V(\nabla\cdot\mathbf{A})\,\mathrm{d}V
      =3\cdot\frac43\pi R^3
      =4\pi R^3.
    \]
    球面上 $\mathbf{A}=R\,\mathbf{e}_r$，外法向也是 $\mathbf{e}_r$，面积元
    $\mathrm{d}S=R^2\sin\theta\,\mathrm{d}\theta\,\mathrm{d}\phi$，故
    $\mathbf{A}\cdot\mathrm{d}\mathbf{S}=R^3\sin\theta\,\mathrm{d}\theta\,\mathrm{d}\phi$。
    \[
      \oint_S\mathbf{A}\cdot\mathrm{d}\mathbf{S}
      =\int_0^{2\pi}\!\mathrm{d}\phi\int_0^{\pi}R^3\sin\theta\,\mathrm{d}\theta
      =2\pi R^3\cdot 2
      =4\pi R^3.
    \]
    两边都是 $4\pi R^3$。

  \item
    \begin{align*}
      (\nabla\times\mathbf{A})_x
      &=\partial_y A_z-\partial_z A_y
       =\partial_y x-\partial_z z=-1,\\
      (\nabla\times\mathbf{A})_y
      &=\partial_z A_x-\partial_x A_z
       =\partial_z y-\partial_x x=-1,\\
      (\nabla\times\mathbf{A})_z
      &=\partial_x A_y-\partial_y A_x
       =\partial_x z-\partial_y y=-1.
    \end{align*}
    所以 $\nabla\times\mathbf{A}=-\mathbf{e}_x-\mathbf{e}_y-\mathbf{e}_z$。
    正方形上 $\mathrm{d}\mathbf{S}=\mathbf{e}_z\,\mathrm{d}x\,\mathrm{d}y$，
    \[
      \int_S(\nabla\times\mathbf{A})\cdot\mathrm{d}\mathbf{S}
      =\int_0^a\int_0^a(-1)\,\mathrm{d}x\,\mathrm{d}y
      =-a^2.
    \]
    法向 $\mathbf{e}_z$ 对应的边界是从上方看的逆时针方向，四条边依次为
    \begin{align*}
      L_1&: y=z=0,\ x:0\to a,
        &&\mathbf{A}\cdot\mathrm{d}\boldsymbol{\ell}=y\,\mathrm{d}x=0,\\
      L_2&: x=a,\ z=0,\ y:0\to a,
        &&\mathbf{A}\cdot\mathrm{d}\boldsymbol{\ell}=z\,\mathrm{d}y=0,\\
      L_3&: y=a,\ z=0,\ x:a\to 0,
        &&\mathbf{A}\cdot\mathrm{d}\boldsymbol{\ell}=y\,\mathrm{d}x
          =\int_a^0 a\,\mathrm{d}x=-a^2,\\
      L_4&: x=z=0,\ y:a\to 0,
        &&\mathbf{A}\cdot\mathrm{d}\boldsymbol{\ell}=z\,\mathrm{d}y=0.
    \end{align*}
    环量是 $-a^2$，与面积分一致。
\end{enumerate}
\end{solution}

%==========================================================================
\begin{problem}
  已知
  $\mathbf{A}\times(\mathbf{B}\times\mathbf{C})
  =(\mathbf{A}\cdot\mathbf{C})\mathbf{B}-(\mathbf{A}\cdot\mathbf{B})\mathbf{C}$.
  \begin{enumerate}
    \item 证明
      \[
        (\mathbf{A}\times\mathbf{B})\times\mathbf{C}
        =(\mathbf{A}\cdot\mathbf{C})\mathbf{B}
         -(\mathbf{B}\cdot\mathbf{C})\mathbf{A}.
      \]
    \item 证明
      \[
        \mathbf{A}\times(\mathbf{B}\times\mathbf{C})
        +\mathbf{B}\times(\mathbf{C}\times\mathbf{A})
        +\mathbf{C}\times(\mathbf{A}\times\mathbf{B})
        =\mathbf{0}.
      \]
    \item 用指标记号证明拉格朗日恒等式
      \[
        (\mathbf{A}\times\mathbf{B})\cdot(\mathbf{C}\times\mathbf{D})
        =(\mathbf{A}\cdot\mathbf{C})(\mathbf{B}\cdot\mathbf{D})
         -(\mathbf{A}\cdot\mathbf{D})(\mathbf{B}\cdot\mathbf{C}).
      \]
  \end{enumerate}
\end{problem}

\begin{solution}
\begin{enumerate}
  \item $(\mathbf{A}\times\mathbf{B})\times\mathbf{C}
    =-\mathbf{C}\times(\mathbf{A}\times\mathbf{B})$。
    把已知公式里的三个矢量换成 $\mathbf{C},\mathbf{A},\mathbf{B}$：
    \[
      \mathbf{C}\times(\mathbf{A}\times\mathbf{B})
      =(\mathbf{C}\cdot\mathbf{B})\mathbf{A}-(\mathbf{C}\cdot\mathbf{A})\mathbf{B}.
    \]
    因此
    \[
      (\mathbf{A}\times\mathbf{B})\times\mathbf{C}
      =(\mathbf{A}\cdot\mathbf{C})\mathbf{B}-(\mathbf{B}\cdot\mathbf{C})\mathbf{A}.
    \]

  \item 已知公式及其轮换给出
    \begin{align*}
      \mathbf{A}\times(\mathbf{B}\times\mathbf{C})
      &=(\mathbf{A}\cdot\mathbf{C})\mathbf{B}-(\mathbf{A}\cdot\mathbf{B})\mathbf{C},\\
      \mathbf{B}\times(\mathbf{C}\times\mathbf{A})
      &=(\mathbf{B}\cdot\mathbf{A})\mathbf{C}-(\mathbf{B}\cdot\mathbf{C})\mathbf{A},\\
      \mathbf{C}\times(\mathbf{A}\times\mathbf{B})
      &=(\mathbf{C}\cdot\mathbf{B})\mathbf{A}-(\mathbf{C}\cdot\mathbf{A})\mathbf{B}.
    \end{align*}
    $\mathbf{A}$ 的系数是 $-(\mathbf{B}\cdot\mathbf{C})+(\mathbf{C}\cdot\mathbf{B})=0$，
    $\mathbf{B}$、$\mathbf{C}$ 的系数同样抵消，三项之和为零。

  \item
    \begin{align*}
      (\mathbf{A}\times\mathbf{B})\cdot(\mathbf{C}\times\mathbf{D})
      &=(\mathbf{A}\times\mathbf{B})_i(\mathbf{C}\times\mathbf{D})_i
       =\varepsilon_{ijk}A_j B_k\,\varepsilon_{ilm}C_l D_m.
    \end{align*}
    用 $\varepsilon_{ijk}\varepsilon_{ilm}=\delta_{jl}\delta_{km}-\delta_{jm}\delta_{kl}$，
    \begin{align*}
      \varepsilon_{ijk}\varepsilon_{ilm}A_j B_k C_l D_m
      &=(\delta_{jl}\delta_{km}-\delta_{jm}\delta_{kl})A_j B_k C_l D_m \\
      &=A_j C_j\,B_k D_k-A_j D_j\,B_k C_k \\
      &=(\mathbf{A}\cdot\mathbf{C})(\mathbf{B}\cdot\mathbf{D})
        -(\mathbf{A}\cdot\mathbf{D})(\mathbf{B}\cdot\mathbf{C}).
    \end{align*}
\end{enumerate}
\end{solution}

\end{document}
